% !TeX root = ../../Book4.tex
% 纲要标题：### 8.4 正合性的障碍
% 纲要位置：计划纲要.md，第 2668 行。

\section{正合性的障碍}
\label{b4:sec:0804}

一次 Ext 或 Tor 的消失能检测投射性、内射性或平坦性，但测试对象必须遍历相应的全部模。整数模与多项式理想的例子可以把正合性的失败具体算出来。


% 纲要内容（仅作写作计划，尚非教材正文）：
%
% * Ext 检验投射性和内射性
% * Tor 检验平坦性
% * 回到第二卷反例解释失败机制
%

\subsection{Ext 与投射性、内射性}
\label{b4:subsec:0804:01}

\begin{theorem}\label{b4:thm:ext-projective-injective}
设 \(R\) 是含幺环，\(P,I\) 是左 \(R\) 模。下列判据成立：
\begin{enumerate}
\item \(P\) 投射，当且仅当对所有左模 \(N\)，\(\operatorname{Ext}_R^1(P,N)=0\)；这也等价于所有正次 \(\operatorname{Ext}_R^n(P,N)\) 消失。
\item \(I\) 内射，当且仅当对所有左模 \(M\)，\(\operatorname{Ext}_R^1(M,I)=0\)；这也等价于所有正次 \(\operatorname{Ext}_R^n(M,I)\) 消失。
\end{enumerate}
\end{theorem}
\begin{proof}
投射或内射时用集中于零次的相应消解，正次 Ext 消失。反过来，若 \(\operatorname{Ext}_R^1(P,-)=0\)，对任意满射 \(B\to C\) 及其核 \(A\)，第二变量长正合列使 \(\operatorname{Hom}_R(P,B)\to\operatorname{Hom}_R(P,C)\) 满，这恰为投射提升性质。若 \(\operatorname{Ext}_R^1(-,I)=0\)，对任意单射 \(A\to B\) 及其余核 \(C\)，第一变量长正合列使 \(\operatorname{Hom}_R(B,I)\to\operatorname{Hom}_R(A,I)\) 满，这恰为内射延拓性质。
\end{proof}

“对所有模”的量词有实质作用。例如 \(\operatorname{Ext}_{\mathbb Z}^1(\mathbb Z/2,\mathbb Q)=0\)，并不能推出 \(\mathbb Z/2\) 投射；这里消失是第二变量内射造成的。又如 \(\operatorname{Ext}_{\mathbb Z}^1(\mathbb Z,\mathbb Z)=0\)，不能推出第二个 \(\mathbb Z\) 内射，因为第一个 \(\mathbb Z\) 已投射。

\subsection{Tor 与平坦性}
\label{b4:subsec:0804:02}

\begin{theorem}\label{b4:thm:tor-flat-criterion}
设 \(R\) 是含幺环，\(N\) 是左 \(R\) 模。下列条件等价：\(N\) 平坦；对每个右 \(R\) 模 \(M\)，\(\operatorname{Tor}_1^R(M,N)=0\)；对每个右模 \(M\) 及每个 \(n>0\)，\(\operatorname{Tor}_n^R(M,N)=0\)。右模的判据以相反侧的左模作测试。
\end{theorem}
\begin{proof}
若 \(N\) 平坦，\(-\otimes_RN\) 正合，故右模投射消解在张量后仍在正次数正合；于是所有正次 Tor 为零。全消失当然蕴含一次消失。若一次均消失，对任意右模单射 \(A\to B\)，令 \(C=B/A\)。长正合列的片段
\[
\operatorname{Tor}_1^R(C,N)\longrightarrow A\otimes_RN
\longrightarrow B\otimes_RN
\]
表明后一个箭头单射。张量积本已右正合，所以它现在正合，即 \(N\) 平坦。反环给出右模版本。
\end{proof}

不必为了判定平坦性计算所有高次 Tor；一次已经足够。相反，只对某一个模计算一次 Tor，通常只能发现反例，不能得出整体平坦性。

\subsection{非正合性的模论反例}
\label{b4:subsec:0804:03}

把 \(0\to\mathbb Z\xrightarrow{\times m}\mathbb Z\to\mathbb Z/m\to0\) 与 \(\mathbb Z/m\)、\(m>1\)，作张量，前一个乘 \(m\) 的映射变为零。长正合列现在告诉我们，丢失的单射性恰由 \(\operatorname{Tor}_1^{\mathbb Z}(\mathbb Z/m,\mathbb Z/m)\cong\mathbb Z/m\) 映入第一个张量项。对同一列施加 \(\operatorname{Hom}_{\mathbb Z}(-,\mathbb Z)\)，乘 \(m\) 的映射不是满射；其余核正是 \(\operatorname{Ext}_{\mathbb Z}^1(\mathbb Z/m,\mathbb Z)\cong\mathbb Z/m\)。第二卷中出现的两种正合性失败由此有了可复用的计量方式。

\ProseParagraphBreak
再取 \(R=k[x,y]\) 的理想 \(\mathfrak m=(x,y)\)。短正合列 \(0\to\mathfrak m\to R\to k\to0\) 与上一节算出的二次 Tor 给出
\[
\operatorname{Tor}_1^R(\mathfrak m,k)\cong
\operatorname{Tor}_2^R(k,k)\cong k.
\]
故 \(\mathfrak m\) 不平坦，尽管它是整环中的无挠子模。这具体说明主理想整环上的无挠平坦结论不能直接推广到任意整环。这里用的是第一变量的长正合列及 \(R\) 的投射性；下一节把这种移位系统写出。
